% ====================================================================
% firmware.tex (design note example)
% ====================================================================

\section{Firmware Interface}
\label{sec:firmware}

The filter is designed around a sample rate of exactly
\qty{10}{\kilo\hertz}, so the firmware has to hold that rate steady. If
conversions were started from the main loop, the interval between them would
vary with whatever else the processor was doing, and that jitter would add
noise to the samples. A hardware timer is used to start each conversion
instead, which means the sample instants are set by the timer and not by the
code.

\subsection{Sampling sequence}

The sequence for each sample is shown in \cref{fig:adcTiming}. The timer
overflows every \qty{100}{\micro\second} and triggers the converter, which
holds the input for its acquisition time and then converts. When the
conversion finishes the converter raises an interrupt, and the interrupt
routine stores the result.

% ====================================================================
% adcTiming.tex
% --------------------------------------------------------------------
% Sampling sequence, drawn with tikz-timing. Two sample periods are
% shown. The widths are illustrative and not to scale.
% ====================================================================

\begin{figure}[htbp]
\centering
\begin{tikztimingtable}[
  timing/slope=0,
  timing/xunit=6mm,
  timing/yunit=6mm
]
  Timer   & 1L 1H 9L 1H 8L \\
  Acquire & 2L 2H 8L 2H 6L \\
  Convert & 4L 4H 6L 4H 2L \\
  IRQ     & 8L 1H 9L 1H 1L \\
\extracode
  \draw[<->, thin] (1,-7.4) -- (11,-7.4)
    node[midway, below, font=\footnotesize] {\qty{100}{\micro\second}};
\end{tikztimingtable}
\caption{The sampling sequence over two sample periods. The timer starts each
conversion, and the interrupt is raised when the result is ready.}
\label{fig:adcTiming}
\end{figure}


The filter's output impedance is low, as it is driven by the op-amp, so the
converter's sample capacitor settles well within the acquisition time and no
extra buffer is needed between the two.

\subsection{Example code}

\Cref{lst:adcSampling} shows the arrangement. The function names stand in for
those of the microcontroller's own hardware library, which differ from one
family to the next, but the structure is the same on any of them.

% ====================================================================
% adcSampling.tex
% --------------------------------------------------------------------
% Timer-triggered ADC sampling. The hardware functions are placeholders
% for those of the target microcontroller's own library.
% ====================================================================

\begin{lstlisting}[caption={Timer-triggered sampling into a ring buffer.}, label={lst:adcSampling}, language=C]
#define SAMPLE_RATE_HZ  10000u
#define BUFFER_LENGTH   256u

static volatile uint16_t buffer[BUFFER_LENGTH];
static volatile uint16_t head;

/* Called once at start-up. The timer starts a conversion every
   100 us, so the sample rate does not depend on the main loop. */
void sampling_init(void)
{
    adc_select_channel(ADC_CHANNEL_SENSOR);
    adc_set_trigger(ADC_TRIGGER_TIMER);
    adc_enable_interrupt();

    timer_set_period_us(1000000u / SAMPLE_RATE_HZ);
    timer_start();
}

/* End-of-conversion interrupt: store the result and move on. */
void adc_isr(void)
{
    buffer[head] = adc_read_result();
    head = (head + 1u) % BUFFER_LENGTH;
}
\end{lstlisting}


The interrupt routine only stores the sample and moves the index on. Anything
slower, such as scaling or further digital filtering, is left to the main
loop, so that the routine is always finished long before the next conversion
completes.
